\section{Related Work}
\label{sec:related}
\textbf{Prompt and program optimisation.} DSPy compiles declarative LM programs by searching instructions and demonstrations against a metric on a training set~\citep{khattab2023dspy}; MIPRO jointly optimises instructions and demonstrations of multi-stage programs~\citep{opsahlong2024mipro}; best-arm-identification formulations select prompts under a limited evaluation budget~\citep{shi2024bai}. These works report the score of the selected configuration on held-out data of a given size; we ask what such a number means when the same items drive selection and scoring, and how it depends on the search budget.

\textbf{Evaluation variability of prompts.} Prompt format alone changes accuracy substantially~\citep{sclar2024formatting}, and prompt-performance distributions can be estimated from few evaluations~\citep{polo2024multiprompt}. This supplies the between-configuration spread that, in our analysis, drives the selection effect.

\textbf{Adaptive data analysis and test-set reuse.} Reusing a holdout adaptively invalidates naive estimates~\citep{dwork2015adaptive}; the Ladder reports thresholded leaderboard scores to limit overfitting~\citep{blum2015ladder}. Empirical studies of test-set reuse find smaller drift than worst-case bounds suggest~\citep{recht2018cifar,mania2019similarity}. We study the same phenomenon for automatic prompt/program search with an explicit selection step, and compare the Ladder-style report as a baseline.

\textbf{Selection-aware inference.} Winner's-curse corrections and conditional/hybrid intervals for selected targets~\citep{andrews2024winners}, simultaneous inference over selected models~\citep{berk2013valid}, Gaussian multiplier-bootstrap bounds for maxima~\citep{chernozhukov2013gaussian}, and the target-of-estimation analysis of cross-validation~\citep{bates2024cv} provide the tools we adapt to reporting a tuned system's score.
